\chapter{Tagging System Index}
\label{chap:Tagging_Index}

\section{Purpose}

The Tagging System provides a multi-ledger classification framework for Engrams.
Each Engram may carry one or more tags describing:
\begin{itemize}
    \item sociological context,
    \item psychological state or identity,
    \item ontological structure,
    \item theological or mythic significance,
    \item epistemic justification,
    \item philosophical framing,
    \item symbolic primitives or glyph-roots,
    \item and PRIME utility functions.
\end{itemize}

Tags enrich Engrams with interpretive metadata, enable structured navigation of
the Golden Engram Library (GEL), and support reconstruction invariants $Inv(E)$.

Tagging is \emph{orthogonal} to W5H but complementary:
W5H provides contextual grounding;  
tags provide conceptual topology.

\section{Ledger Overview}

The Tagging System is divided into eight ledgers, each serving a distinct
interpretive function.

\subsection{SOCI Ledger (Sociology)}
Focuses on:
\begin{itemize}
    \item social relationships and role dynamics,
    \item group structures and collective behaviors,
    \item cultural frames and societal meaning-making.
\end{itemize}

Typical tags:
\begin{quote}
\texttt{soci:ritual}, \texttt{soci:collective-action}, \texttt{soci:identity-group}
\end{quote}

\subsection{PSY Ledger (Psychology)}
Encodes:
\begin{itemize}
    \item affective states,
    \item personality structures,
    \item motive and disposition analysis.
\end{itemize}

Typical tags:
\begin{quote}
\texttt{psy:arousal}, \texttt{psy:attachment}, \texttt{psy:defense-mechanism}
\end{quote}

\subsection{ONTO Ledger (Ontology)}
Tracks:
\begin{itemize}
    \item metaphysical categories,
    \item structural relationships,
    \item entity classifications.
\end{itemize}

Typical tags:
\begin{quote}
\texttt{onto:entity}, \texttt{onto:relation}, \texttt{onto:substance}
\end{quote}

\subsection{THEO Ledger (Theology)}
Captures:
\begin{itemize}
    \item symbolic or mythic resonance,
    \item religious context,
    \item templating through sacred narratives.
\end{itemize}

Typical tags:
\begin{quote}
\texttt{theo:ritual}, \texttt{theo:myth}, \texttt{theo:archetype}
\end{quote}

\subsection{EPI Ledger (Epistemology)}
Covers:
\begin{itemize}
    \item justification structures,
    \item knowledge claims,
    \item evidential or inferential status.
\end{itemize}

Typical tags:
\begin{quote}
\texttt{epi:axiom}, \texttt{epi:empirical}, \texttt{epi:inference}
\end{quote}

\subsection{PHIL Ledger (Philosophy)}
Describes:
\begin{itemize}
    \item logical, ethical, or metaphysical framing,
    \item philosophical schools and stances.
\end{itemize}

Typical tags:
\begin{quote}
\texttt{phil:virtue}, \texttt{phil:logic}, \texttt{phil:phenomenology}
\end{quote}

\subsection{GEN Ledger (Generic Symbolic Primitives)}
Contains:
\begin{itemize}
    \item fundamental symbolic primitives,
    \item glyph roots,
    \item universal structures or motifs.
\end{itemize}

Typical tags:
\begin{quote}
\texttt{gen:root}, \texttt{gen:glyph}, \texttt{gen:primitive}
\end{quote}

\subsection{PRIME Ledger (Utility Tags)}
Utility-specific tags used for:
\begin{itemize}
    \item inline operations,
    \item ribbons or anchors,
    \item index management,
    \item structural metadata.
\end{itemize}

Typical tags:
\begin{quote}
\texttt{prime:anchor}, \texttt{prime:important}, \texttt{prime:role}
\end{quote}

\section{Tag Schema}

A Tag $T$ is defined as:

\begin{quote}
\[
T =
\langle ledger,\, key,\, gloss,\, class\_bindings,\, constraints \rangle
\]
\end{quote}

Where:
\begin{itemize}
    \item \textbf{ledger}: one of \{SOCI, PSY, ONTO, THEO, EPI, PHIL, GEN, PRIME\}
    \item \textbf{key}: string identifier under ledger
    \item \textbf{gloss}: natural-language description
    \item \textbf{class\_bindings}: optional (C/L/M/S/O/P)
    \item \textbf{constraints}: rules for when this tag may be applied
\end{itemize}

Example:
\begin{quote}
\texttt{onto:relation}  
Gloss: ``Describes a structural connection between two entities.''  
Bindings: S-class preferred.  
\end{quote}

\section{Tag Application Rules}

\subsection{Tagging Requirements}

\begin{enumerate}
    \item All W5H-Applied Engrams MUST contain at least one tag.
    \item Tags MUST be semantically compatible with:
    \begin{itemize}
        \item Engram type (Mote, Constructor, W5H-applied),
        \item W5H fields,
        \item ZED Core identity constraints.
    \end{itemize}
    \item A tag MUST NOT contradict:
    \begin{itemize}
        \item governing P-class principles,
        \item canonical C-chain history,
        \item reconstruction invariants $Inv(E)$.
    \end{itemize}
\end{enumerate}

\subsection{Class Binding Rules}

Tags may carry C/L/M/S/O/P class bindings.

\begin{itemize}
    \item C/P-bound tags require HITL review.
    \item S-class tags affect structural interpretation.
    \item O-class tags interact with narrative branches.
    \item L-class tags are ephemeral and may evaporate.
\end{itemize}

\section{Tag Interactions with GEL}

Tags serve as primary indices for the Golden Engram Library (GEL):

\begin{itemize}
    \item Tag → Subject Lattice mapping
    \item Tag co-occurrence analysis
    \item Ontological adjacency
    \item Thematic clustering
\end{itemize}

GEL uses tags to:
\begin{itemize}
    \item cross-reference Engrams,
    \item build conceptual neighborhoods,
    \item support efficient reconstruction,
    \item maintain harmomic coherence across branches.
\end{itemize}

\section{Tagging Workflow}

For each Engram $E$:
\begin{enumerate}
    \item Collect W5H metadata.
    \item Draft tag set $\{T_i\}$.
    \item Validate each tag against:
    \begin{itemize}
        \item ledger constraints,
        \item W5H consistency,
        \item ZED/K constraints,
        \item reconstruction invariants.
    \end{itemize}
    \item Submit P-class or C-class tags for HITL approval.
    \item Commit Engram to:
    \begin{itemize}
        \item \texttt{Tag\_Applications/}
        \item GEL index
        \item Subject Lattice
    \end{itemize}
\end{enumerate}

\section{Summary}

The Tagging System provides the conceptual topology through which Engrams can be:
\begin{itemize}
    \item classified,
    \item related,
    \item searched,
    \item reconstructed,
    \item and meaningfully interpreted.
\end{itemize}

Where W5H answers the contextual ``who / what / when / where / why / how,''  
the tagging system answers the structural ``in what conceptual universe does this Engram live?''

Together, they form the backbone of the CME's semantic and ontological stability.
